\section{Impl\'ementation}
\label{sec:impl}

\subsection{Choix techniques}

Les choix techniques sont r\'esum\'es dans le tableau~\ref{tab:tech}. Tous les param\`etres
num\'eriques proviennent de \texttt{config.py}.

\begin{table}[H]
  \centering
  \caption{Choix techniques et justifications.}
  \label{tab:tech}
  \small
  \begin{tabular}{@{}p{2.4cm}p{3.2cm}p{7.8cm}@{}}
    \toprule
    \textbf{Composant} & \textbf{Choix} & \textbf{Justification} \\
    \midrule
    Langage & Python 3.13 & \'Ecosyst\`eme NLP mature (sentence-transformers, ChromaDB, Streamlit). \\
    Embeddings & all-MiniLM-L6-v2 & Mod\`ele compact ($\sim$80~Mo), rapide en CPU, adapt\'e \`a la recherche s\'emantique. \\
    Base vectorielle & ChromaDB (persistant) & Int\'egration simple, pas de serveur externe, persistance locale. \\
    LLM & Groq --- llama-3.3-70b-versatile & API rapide, mod\`ele open-weight performant. \\
    Interface & Streamlit + CLI & CLI pour le d\'eveloppement ; Streamlit pour la d\'emo interactive. \\
    D\'ecoupage & 512 / overlap 50 & Compromis granularit\'e / contexte par chunk. \\
    Top-k & 3 & Suffisant pour couvrir la r\'eponse sans diluer le contexte. \\
    \bottomrule
  \end{tabular}
\end{table}

\subsection{Pipeline de r\'ecup\'eration et de g\'en\'eration}

Le module \texttt{src/rag\_pipeline.py} impl\'emente les fonctions centrales :
\begin{itemize}[noitemsep]
  \item \texttt{retrieve()} : encode la question et interroge ChromaDB ;
  \item \texttt{build\_prompt()} : assemble le contexte r\'ecup\'er\'e et la question ;
  \item \texttt{generate\_answer()} : appelle le LLM Groq et retourne la r\'eponse ainsi que
        les chunks sources.
\end{itemize}

Le template de prompt demande explicitement au mod\`ele de r\'epondre \textbf{uniquement} \`a partir
du contexte fourni et d'indiquer s'il ne conna\^it pas la r\'eponse, afin de limiter les
hallucinations. Un extrait est donn\'e en annexe~\ref{ann:code}.

\subsection{Interface utilisateur}

\subsubsection{Ligne de commande (op\'erationnelle)}

La commande \texttt{python -m src.rag\_pipeline} lance une boucle interactive : l'utilisateur
pose une question, le syst\`eme affiche la r\'eponse et les documents sources. Cette interface
a servi de base pour valider le pipeline avant l'\'evaluation automatique.

\subsubsection{Streamlit (impl\'ement\'ee, validation en cours)}

L'application \texttt{app.py} organise l'interface en trois onglets :
\begin{itemize}[noitemsep]
  \item \textbf{Chat} : conversation avec affichage du temps de r\'eponse, des chunks utilis\'es,
        d'un score de fid\'elit\'e et des extraits sources ;
  \item \textbf{\'Evaluation} : chargement d'un jeu de questions (JSON/CSV), lancement de
        l'\'evaluation, m\'etriques agr\'eg\'ees, graphiques Plotly et export CSV/PDF ;
  \item \textbf{Gestion du corpus} : statistiques, ajout de documents (PDF/TXT), s\'election des
        sources et reconstruction de l'index vectoriel.
\end{itemize}

La barre lat\'erale permet de choisir le mod\`ele LLM Groq, d'ajuster \texttt{top-k} et le seuil
de similarit\'e. Les tests finaux de cette interface et la pr\'eparation de la d\'emonstration
orale constituent le travail restant identifi\'e en section~\ref{sec:etat}.

\subsection{Scripts principaux}

\begin{table}[H]
  \centering
  \caption{Scripts du projet et r\^ole.}
  \label{tab:scripts}
  \small
  \begin{tabular}{@{}p{4.5cm}p{9cm}@{}}
    \toprule
    \textbf{Fichier} & \textbf{R\^ole} \\
    \midrule
    \texttt{collect\_corpus.py} & T\'el\'echargement des articles Wikipedia \\
    \texttt{src/prepare\_data.py} & Nettoyage et chunking \\
    \texttt{src/build\_index.py} & Embeddings et index ChromaDB \\
    \texttt{src/rag\_pipeline.py} & R\'ecup\'eration + g\'en\'eration (CLI) \\
    \texttt{app.py} & Interface Streamlit \\
    \texttt{evaluation/evaluate.py} & \'Evaluation batch sur le jeu de test \\
    \texttt{visualizations/plot\_results.py} & G\'en\'eration des graphiques \\
    \bottomrule
  \end{tabular}
\end{table}
